\section{Stabilizers and asymptotics}\label{sec:types}

\begin{theorem}\label{thm:stabilizers}
For every $n\ge0$, the exact stabilizers of three-row tilings
are distributed as in Table \ref{tab:stabilizers}.
\end{theorem}

\begin{table}[htbp]
\centering
\caption{Tilings and orbits classified by exact stabilizer.}
\label{tab:stabilizers}
\begin{tabular}{cccc}
\toprule
Stabilizer & Tilings & Orbit size & Orbits\\
\midrule
$G$ & $1$ & $1$ & $1$\\
$\{e,h\}$ & $0$ & $2$ & $0$\\
$\{e,v\}$ & $v_n-1$ & $2$ & $(v_n-1)/2$\\
$\{e,\rho\}$ & $r_n-1$ & $2$ & $(r_n-1)/2$\\
$\{e\}$ & $d_n-v_n-r_n+1$ & $4$ & $(d_n-v_n-r_n+1)/4$\\
\bottomrule
\end{tabular}
\end{table}

\begin{proof}
By Lemma \ref{lem:h-odd}, the only tiling fixed by $h$ is the
tiling in which each of the three rows is tiled horizontally.
It is also fixed by $v$ and $\rho$.
A tiling fixed by any two distinct nonidentity elements is
fixed by their product, and hence by all of $G$.
Thus every pairwise intersection, as well as the triple
intersection, of the three nonidentity fixed sets is this
single tiling.

There are consequently $v_n-1$ tilings with exact stabilizer
$\{e,v\}$, $r_n-1$ with exact stabilizer $\{e,\rho\}$, none with
exact stabilizer $\{e,h\}$, and $d_n-v_n-r_n+1$ with trivial
stabilizer. Since $G$ is abelian, the stabilizer is constant
on an orbit. Dividing by the corresponding orbit size gives
the table. These five subgroups exhaust all subgroups of $G$.
\end{proof}

For example, when $n=8$ the fixed-point counts are
$d_8=29681$, $v_8=265$, and $r_8=153$.
Among the $7525$ orbits, $7316$ have trivial stabilizer,
$132$ have stabilizer $\{e,v\}$, $76$ have stabilizer
$\{e,\rho\}$, and one has stabilizer $G$.

\begin{theorem}\label{thm:asym}
Let $\lambda=2+\sqrt3$. If a tiling is chosen uniformly from
$\mathcal T_n$, let $p_n$ be the probability that it has a
nonidentity symmetry. Then, as $j\to\infty$,
\begin{equation}\label{eq:prob-asym}
 p_{2j}\sim(\lambda-1)\lambda^{-j},\qquad
 p_{2j+1}\sim\frac{\lambda+1}{\lambda}\lambda^{-j}.
\end{equation}
In particular, $p_n=O(\lambda^{-n/2})\to0$, and
\begin{equation}\label{eq:orbit-asym}
 a(n)=\frac{d_n}{4}+O(\lambda^{n/2})
 \sim\frac{3+\sqrt3}{24}\lambda^n.
\end{equation}
\end{theorem}
\begin{proof}
Solving \eqref{eq:d-rec} with its initial values gives
\begin{equation}\label{eq:d-closed}
 d_n=c_+\lambda^n+c_-\lambda^{-n},\qquad
 c_\pm=\frac{3\pm\sqrt3}{6}.
\end{equation}
By Table \ref{tab:stabilizers},
$p_n=(v_n+r_n-1)/d_n$, whereas \eqref{eq:parity-fixed} gives
\[
 v_{2j}+r_{2j}=d_{j+1}-d_j,\qquad
 v_{2j+1}+r_{2j+1}=d_{j+1}+d_j.
\]
The leading terms of these numerators are
$c_+(\lambda-1)\lambda^j$ and $c_+(\lambda+1)\lambda^j$,
respectively. The corresponding denominators have leading
terms $c_+\lambda^{2j}$ and $c_+\lambda^{2j+1}$.
This proves \eqref{eq:prob-asym}. Finally,
$a(n)=(d_n+1+v_n+r_n)/4$ gives \eqref{eq:orbit-asym}.
\end{proof}

Uniform sampling of orbits gives a different probability
space. The number of orbits with nontrivial stabilizer is
$(v_n+r_n)/2$, so their proportion is
$(v_n+r_n)/(2a(n))$, which also tends to zero.
